\chapter[Entanglement equilibrium]{Entanglement Equilibrium and the Modular Hamiltonian}
\label{ch:entanglement-equilibrium}

\paragraph{Learning outcomes.}
After this chapter the reader should be able to:
\begin{enumerate}
  \item write the reduced state \(\rho_D\), the modular Hamiltonian
        \(H_D\), and the von Neumann entropy \(S_D\) of a causal
        diamond in the matrix condensate;
  \item state Jacobson's entanglement-equilibrium condition
        \(\delta S_D=0\) and distinguish it from the Einstein
        equation it will later imply;
  \item obtain the first variation \(\delta S_D=\delta\langle H_D\rangle\)
        at fixed diamond volume, including the meaning of that
        constraint;
  \item define the mutual-information potential \(U_{\mathrm{ent}}\)
        used for two condensed cores;
  \item enter Chapter~7 with a Ward identity ready to be applied to
        \(T_{\mu\nu}\).
\end{enumerate}

Chapter~5 named the candidate graviton \(T_{\mu\nu}\) and left it
with a generic mass \(m_T\). The mass is not removed by a further
term in \(V(\Phi^\dagger\Phi)\). It is removed, if at all, by a
stationarity condition on entanglement. This chapter writes that
condition in a form that can be varied. The variation itself is
Chapter~7.

\section{The logical role of this chapter}
\label{sec:ch6-role}

There are two different claims that are easy to conflate.

\begin{enumerate}
  \item \textbf{Jacobson's theorem} (in the form used here).
        For a quantum field on a prescribed geometry, the
        first-order stationarity of vacuum entanglement in small
        geodesic balls, at fixed volume, is equivalent to the
        semiclassical Einstein equation. This direction assumes a
        metric and a quantum field, and obtains \(G_{\mu\nu}\propto
        \langle T_{\mu\nu}\rangle\).
  \item \textbf{The NAMCG claim.}
        The microscopic condensate state is itself an entanglement-
        stationary state on the diamonds of its \emph{emergent}
        metric. Stationarity is therefore a property of \(\rho\),
        not an extra thermodynamic axiom pasted onto a metric theory.
        Applied to the tensor block of Chapter~5, it forbids
        \(m_T^2\neq 0\).
\end{enumerate}

Claim~1 is a theorem in the literature, with stated hypotheses
(conformal matter at first order; a conjecture for non-conformal
matter). Claim~2 is the load-bearing physical hypothesis of NAMCG.
This chapter formulates Claim~2 so that it can be used. It does not
pretend that Claim~2 has been derived from the path integral of
Chapter~3 with the same status as Claim~1.

\begin{assumption}[Microscopic entanglement equilibrium]
  \label{ass:EE}
  Let \(\rho\) be the vacuum state of the matrix condensate in the
  condensed phase of Chapters~3--4. For every sufficiently small
  causal diamond \(D\) constructed from the emergent metric
  \(g_{\mu\nu}\),
  \begin{equation}
    \label{eq:EE}
    \delta S_D=0
  \end{equation}
  at fixed spatial volume of the maximal slice of \(D\), to first
  order in variations of the condensate that preserve the local
  vacuum character of \(\rho\).
\end{assumption}

The remainder of the chapter unpacks every symbol in that statement.

\section{Diamonds in the emergent geometry}
\label{sec:ch6-diamonds}

The causal diamond \(D(p,q)=J^+(p)\cap J^-(q)\) was defined in
\eqref{eq:diamond}. From here on the causal structure is that of the
\emph{emergent} metric~\eqref{eq:g-from-e}, not of the auxiliary
scaffold \(g^{(0)}\). For a linearized analysis about Minkowski
spacetime the distinction is invisible at zeroth order and appears
at first order through \(h_{\mu\nu}\).

A small diamond of radius \(\ell\) is one for which
\(\ell^2\|R_{\mu\nu\rho\sigma}\|\ll 1\). In Riemann normal
coordinates~\eqref{eq:RNC} its maximal-slice three-volume admits the
expansion
\begin{equation}
  \label{eq:vol-exp}
  V(D)
  =\frac43\pi\ell^3
   -\frac{4\pi\ell^5}{15}\,
     \bigl(\tfrac13 R+\text{anisotropic Ricci terms}\bigr)
   +O(\ell^7).
\end{equation}
``Fixed volume'' in Assumption~\ref{ass:EE} means that the first
variation is taken subject to \(\delta V(D)=0\). The geometric
meaning of that constraint is that a conformal breathing of the
diamond is not counted as an independent entanglement variation; it
has been used to fix the comparison background. This is the same
gauge-like restriction Jacobson imposes, and it is why the conformal
mode \(\psi\) of Chapter~5 and the tensor mode \(T_{\mu\nu}\) enter
the first law differently.

Dyadic labels \(D_{j,n}\) are a nested bookkeeping device. Nothing
physical depends on bisection. Every identity written for a generic
small \(D\) applies to each dyadic diamond.

\section{Reduced states and the modular Hamiltonian}
\label{sec:ch6-modular}

Let \(\mathcal{A}_D\) be the algebra of condensate observables
supported in \(D\). In Presentation~B these are functions of
\(\Phi(x)\) and \(A_\mu(x)\) for \(x\in D\). In Presentation~A they
are the operators generated by the spectral projections of
\(\hat X^\mu\) whose eigenvalues lie in \(D\), after the dictionary
of Assumption~\ref{ass:dictionary}.

\begin{definition}[Diamond state and modular Hamiltonian]
  \label{def:rhoD}
  Given a global state \(\rho\) of the condensate,
  \begin{equation}
    \label{eq:rhoD}
    \rho_D=\mathrm{Tr}_{\mathcal{A}_D'}\rho,
  \end{equation}
  where \(\mathcal{A}_D'\) is the commutant (the complement algebra).
  The modular Hamiltonian and von Neumann entropy are
  \begin{equation}
    \label{eq:HD-SD}
    H_D=-\ln\rho_D,
    \qquad
    S_D=-\mathrm{Tr}(\rho_D\ln\rho_D)=\langle H_D\rangle_{\rho_D}.
  \end{equation}
\end{definition}

For a free field in Rindler space, the Bisognano--Wichmann theorem
identifies \(H_D\) with \(2\pi/\kappa\) times the boost generator
that fixes the edge \(\mathcal{B}=\partial D\cap\Sigma\). For a small
diamond in a slowly varying background the same identification holds
to leading order in \(\ell\):
\begin{equation}
  \label{eq:HD-boost}
  H_D
  =\frac{2\pi}{\kappa}
   \int_{\Sigma\cap D}
   \mathrm{d}\Sigma^\mu\,\xi^\nu T_{\mu\nu}^{\mathrm{op}}
   +H_{\mathrm{area}}
   +O(\ell^3).
\end{equation}
Here \(\xi^\mu\) is the approximate conformal Killing vector of \(D\),
\(\kappa\) is its surface gravity at \(\mathcal{B}\),
\(T_{\mu\nu}^{\mathrm{op}}\) is the microscopic stress tensor of the
condensate fluctuations, and \(H_{\mathrm{area}}\) is the area term
that will later match \(A/(4G_{\mathrm{eff}})\).

Equation~\eqref{eq:HD-boost} is the working formula of this book. It
is exact only for free fields in Rindler space. For interacting
matrix fluctuations it is the first term in an operator product
expansion on the diamond. Higher terms are absorbed into the
\(O(\ell^3)\) remainder and into the definition of \(H_{\mathrm{area}}\).

\section{Replica formula}
\label{sec:ch6-replica}

When a direct construction of \(\rho_D\) is unavailable, the replica
trick supplies \(S_D\). For integer \(n\ge 1\),
\begin{equation}
  \label{eq:replica}
  S_D
  =\lim_{n\to 1}
   \bigl(1-n\partial_n\bigr)
   \ln\mathrm{Tr}\,(\rho_D)^n,
\end{equation}
and \(\mathrm{Tr}\,(\rho_D)^n\) is computed as a path integral of
the microscopic action~\eqref{eq:Smicro} on an \(n\)-sheeted cover
of \(M\) branched over the edge \(\mathcal{B}\). The first variation
\(\delta S_D\) then becomes a variation of that branched partition
function with respect to the background frame.
In the Euclidean replica formulation, after Wick rotation, the path
integral uses
\(S_E=\int\mathrm{d}^4x\,\sqrt{g}\,\mathrm{Tr}\bigl(
g^{\mu\nu}(D_\mu\Phi)^\dagger D_\nu\Phi+V\bigr)\) as the continuation of
\eqref{eq:Smicro}. Matching the corrected Lorentzian sign of
\eqref{eq:Smicro} to Presentation~A remains unproved under
Assumption~\ref{ass:dictionary}.

Two remarks are enough for Chapter~7.

\begin{itemize}
  \item The replica construction is well defined in Presentation~B,
        where a path integral exists. In Presentation~A it requires
        the dictionary of Assumption~\ref{ass:dictionary}.
  \item Differentiating~\eqref{eq:replica} at \(n=1\) produces an
        insertion of the modular Hamiltonian. The first variation
        \(\delta S_D=\delta\langle H_D\rangle\) used below is therefore
        not a second formula; it is the replica formula at first
        order.
\end{itemize}

\section{The first law of the diamond}
\label{sec:ch6-first-law}

Vary Definition~\ref{def:rhoD} at fixed diamond, then allow the
diamond itself to vary with the metric. The result, to first order
and at fixed volume, is the entanglement first law
\begin{equation}
  \label{eq:first-law}
  \delta S_D
  =\delta\langle H_D\rangle
  =\frac{2\pi}{\kappa}
   \int_{\Sigma\cap D}
   \mathrm{d}\Sigma^\mu\,\xi^\nu
   \delta\langle T_{\mu\nu}^{\mathrm{op}}\rangle
   +\frac{\delta A}{4\Geff}
   +O(\ell^3).
\end{equation}
The area variation \(\delta A\) is computed from the emergent metric.
The identification of its prefactor with \(1/(4\Geff)\) is not
justified in this chapter; it is the Bekenstein--Hawking relation
that Chapter~8 will match to the heat kernel. Until then,
\(1/(4\Geff)\) is a name for the coefficient of \(\delta A\).

Assumption~\ref{ass:EE} now reads
\begin{equation}
  \label{eq:first-law-zero}
  \frac{2\pi}{\kappa}
   \int_{\Sigma\cap D}
   \mathrm{d}\Sigma^\mu\,\xi^\nu
   \delta\langle T_{\mu\nu}^{\mathrm{op}}\rangle
   +\frac{\delta A}{4\Geff}
  =0,
\end{equation}
up to \(O(\ell^3)\), for every small diamond and every allowed
variation of the condensate.

If \(\delta\langle T_{\mu\nu}^{\mathrm{op}}\rangle\) were the
stress tensor of ordinary matter on a prescribed geometry, and if
\(\delta A\) were computed from a prescribed Einstein metric,
\eqref{eq:first-law-zero} would be Jacobson's derivation of
\(G_{\mu\nu}=8\pi\Geff\,\langle T_{\mu\nu}\rangle\). In NAMCG both
\(\delta\langle T^{\mathrm{op}}\rangle\) and \(\delta A\) are
functionals of the \emph{same} fluctuation \(\delta\Phi\). That is
the leverage used in Chapter~7.

\section{Allowed variations}
\label{sec:ch6-allowed}

Not every deformation of \(\Phi\) is an allowed variation in
Assumption~\ref{ass:EE}. The hypothesis is a statement about
perturbations that remain locally vacuum-like: they do not inject
an arbitrary excited-state entropy into \(D\).

\begin{definition}[Allowed variation]
  \label{def:allowed}
  A first-order deformation \(\delta\Phi\) is \emph{allowed} if
  \begin{enumerate}
    \item it is smooth on a neighbourhood of \(D\);
    \item it preserves the polar vacuum to zeroth order outside a
          compact region, so that \(\rho\) remains a condensed-phase
          vacuum at infinity;
    \item its restriction to each algebraic block of
          Chapter~5 may be chosen independently, after gauge fixing
          of diffeomorphisms and local Lorentz transformations.
  \end{enumerate}
\end{definition}

Condition~3 is the reason the tensor, vector, and conformal blocks
can be tested separately. If a variation supported only on
\(T_{\mu\nu}\) violates~\eqref{eq:first-law-zero} unless
\(m_T=0\), that is a constraint on the tensor block alone.

\section{Mutual information of two cores}
\label{sec:ch6-MI}

The same entropy data define an interaction between two disjoint
condensed regions \(A\) and \(B\) (two compact-object cores, two
stellar interiors, two laboratory samples).

\begin{definition}[Entropic potential]
  \label{def:Uent}
  \begin{equation}
    \label{eq:Uent-def}
    U_{\mathrm{ent}}(A,B)
    =-\lambda I(A:B)
    =-\lambda\bigl[S(A)+S(B)-S(A\cup B)\bigr],
  \end{equation}
  with \(\lambda>0\) a microscopic coupling of dimension energy.
\end{definition}

Because \(I(A:B)\ge 0\), one has \(U_{\mathrm{ent}}\le 0\). Increasing
the mutual information lowers the potential: the effective force is
attractive. For two small balls whose proper separation \(d\) is
large compared with their radii but small compared with the curvature
radius,
\begin{equation}
  \label{eq:I-falloff}
  I(A:B)
  \;\sim\;
  c_{p_*}
  \left(\frac{\ell_A\ell_B}{d^2}\right)^{p_*}
  \qquad
  (d\gg \ell_A,\ell_B).
\end{equation}
Here \(p_*\) is the smallest sum of scaling dimensions of operators
whose product has vacuum quantum numbers, and \(c_{p_*}\) is a
dimensionless, theory- and shape-dependent coefficient. For the
exchange of a pair of operators of dimension \(\Delta_*\),
\(p_*=2\Delta_*\); other channels can set a different leading power.
Cardy~\cite{Cardy2013MI} derives this scaling for the R\'enyi mutual
information of compact regions in the ground state of a
flat-space conformal field theory, and gives the von Neumann limit
explicitly for spherical regions in the free-scalar case in
\(2+1\) and \(3+1\) dimensions, where for \(3+1\) dimensions
\(p_*=2\), i.e.\ a tail proportional to \(\ell_A^2\ell_B^2/d^4\).
Applying~\eqref{eq:I-falloff} to the von Neumann mutual information
of other theories assumes that the \(n\to1\) continuation holds, and
applying it to NAMCG cores, which need not be vacuum balls of a
conformal theory, is a further assumption. Neither \(p_*\) nor
\(c_{p_*}\) is computed here for NAMCG. The static \(U_{\mathrm{ent}}\)
defines a conservative interaction. Its separation dependence alone
does not determine a velocity-dependent dissipative force or a
retarded response. The Newtonian laboratory of Chapter~15 adopts
\begin{equation}
  \label{eq:airrev-preview}
  \mathbf{a}_{\mathrm{irrev}}
  =-\eta\,\frac{1}{d^4}\,\mathbf{v}
\end{equation}
as a phenomenological working law. Equation~\eqref{eq:airrev-preview}
is not derived from~\eqref{eq:I-falloff}; its microscopic dynamical
origin is Problem~\ref{sec:status-I}. It is flagged here so that
Chapter~11 knows which entropy functional it must vary covariantly.

The conservation question is already visible. Mutual information is
built from scalars \(S(\,\cdot\,)\). A covariant variation of
\(U_{\mathrm{ent}}\) with respect to the metric defines a tensor
\(\Tent_{\mu\nu}\) by the same rule as~\eqref{eq:T-def}. Whether
\(\nabla^\mu\Tent_{\mu\nu}=0\) holds on shell is not automatic from
the sign of \(\lambda\). It is a calculation, reserved for Chapter~11.

\section{What stationarity does \emph{not} say}
\label{sec:ch6-does-not}

Assumption~\ref{ass:EE} does not assert any of the following.

\begin{itemize}
  \item That \(S_D\) is a maximum rather than a saddle. Jacobson's
        original language is maximization at fixed volume; the
        identities used below need only first-order stationarity.
  \item That the area term in~\eqref{eq:first-law} has already been
        matched to \(\Geff\). That match is Chapter~8.
  \item That \(m_T=0\). Stationarity is a condition on variations.
        Whether the tensor kinetic operator is compatible with that
        condition is Chapter~7.
  \item That \(\Lambda_{\mathrm{eff}}\) is small. Vacuum entropy of
        the diamond includes a piece that will source a
        cosmological term. Chapter~12 treats it.
  \item That the replica path integral has been computed at two
        loops. The book uses the first law~\eqref{eq:first-law} as
        the working identity, not a finished replica evaluation of
        \(S_D\) for the matrix potential~\eqref{eq:V}.
\end{itemize}

\section{The Ward identity that Chapter~7 will use}
\label{sec:ch6-ward}

Collecting~\eqref{eq:first-law-zero}, Definition~\ref{def:allowed},
and the split~\eqref{eq:algebraic-split}, the usable statement is:

\begin{proposition}[Entanglement Ward identity]
  \label{prop:ward}
  Under Assumptions~\ref{ass:EE} and~\ref{ass:dictionary}, every
  allowed variation supported on a single algebraic block
  \(\psi\), \(A_\mu\), or \(T_{\mu\nu}\) satisfies
  \begin{equation}
    \label{eq:ward}
    \int_{\Sigma\cap D}
    \mathrm{d}\Sigma^\mu\,\xi^\nu
    \delta\langle T_{\mu\nu}^{\mathrm{op}}\rangle
    +\frac{\kappa}{8\pi\Geff}\,\delta A
    =0
  \end{equation}
  to first order in \(\ell\), at fixed \(V(D)\). The operator
  \(T_{\mu\nu}^{\mathrm{op}}\) is the metric variation of the
  quadratic action of that block.
\end{proposition}

The proof of Proposition~\ref{prop:ward} is the chain of this
chapter: Definition~\ref{def:rhoD} to the first law
\eqref{eq:first-law} to Assumption~\ref{ass:EE}, restricted to one
block by Definition~\ref{def:allowed}. No property of \(T_{\mu\nu}\)
beyond its being one algebraic block has been used. That is
intentional. Chapter~7 substitutes the quadratic action
\eqref{eq:ST-generic} into~\eqref{eq:ward} and reads off \(m_T=0\).

\section{Exercises}
\label{sec:ch6-exercises}

\begin{enumerate}
  \item \textbf{Volume constraint.}
    From~\eqref{eq:vol-exp}, show that a pure conformal variation
    \(\delta g_{\mu\nu}=\chi\eta_{\mu\nu}\) changes \(V(D)\) at
    order \(\ell^3\). Why must this variation be excluded, or
    compensated, if stationarity is imposed at fixed volume?
  \item \textbf{Rindler check.}
    In Rindler space with Unruh temperature \(T=\kappa/2\pi\),
    verify that \(\delta S=\delta E/T\) for a boost-energy
    perturbation is the restriction of~\eqref{eq:first-law} to a
    diamond that has become a Rindler wedge.
  \item \textbf{Replica derivative.}
    Starting from~\eqref{eq:replica}, show that
    \(\partial_n\ln\mathrm{Tr}\,\rho^n|_{n=1}=\langle\ln\rho\rangle\),
    and hence that \(S=\langle H\rangle\) at \(n=1\).
  \item \textbf{Allowed versus thermal variations.}
    Give one example of a \(\delta\Phi\) that changes the particle
    content of \(D\) (a one-particle excitation) and explain why
    Assumption~\ref{ass:EE} is silent about it.
  \item \textbf{Mutual information positivity.}
    Prove \(I(A:B)\ge 0\) from relative entropy. Deduce
    \(U_{\mathrm{ent}}\le 0\) for \(\lambda>0\).
  \item \textbf{Static information versus drag.}
    Explain why a static power law for \(I(A:B)\) does not determine
    the low-frequency retarded force response or its drag exponent.
    What dynamical information is required to calculate that
    response?
  \item \textbf{Two metrics.}
    Write one paragraph explaining why the diamond of this chapter
    is built from \(g_{\mu\nu}\) rather than from \(g^{(0)}_{\mu\nu}\).
    What would go wrong, conceptually, if stationarity were imposed
    on auxiliary-scaffold diamonds instead?
\end{enumerate}

\section{Bridge to Chapter~7}
\label{sec:ch6-bridge}

The candidate graviton \(T_{\mu\nu}\) has a quadratic action
\eqref{eq:ST-generic} with an unknown mass \(m_T\). The condensate
vacuum is assumed to obey the entanglement Ward identity
\eqref{eq:ward} on every small diamond of the emergent geometry.
Chapter~7 substitutes a compactly supported tensor variation into
that identity. The area term tracks the Einstein operator. The
stress-tensor term tracks \(m_T^2 T_{\mu\nu}\). Compatibility of the
two forces \(m_T=0\), after which the surviving kinetic operator is
the linearized Einstein operator and the infrared spectrum contains
two massless TT polarizations---provided the conformal and vector
blocks are separately gapped or eaten, as already flagged.